\section{Definitions}\label{sec:agg-creds-definitions}

\subsection{Aggregate anonymous credentials}

An aggregate anonymous credential scheme involves four types of participants.
The registration authority is responsible for registering users. Issuers certify
attributes. For simplicity we assume that each issuer is responsible for a
single attribute type, and is the only issuer for that attribute. Users each
register with the registration authority, obtain credentials from the relevant
issuers for the attribute values they hold, and then use these credentials to
make presentations. A presentation is a user provably disclosing the values of a
subset of their attributes. Verifiers check users' claims. We write $\att{i}$ for
the attribute type certified by issuer $i$, and $\atv{i}$ for the value of that
attribute held by a user. Issuer $i$ has key pair $(\isk{i}, \ipk{i})$, and the
registration authority $(\rask, \rapk)$.

We assume the registration authority is honest but curious. It follows the
protocol and keeps its signing key secret, but may share the rest of its view
with other participants to try to identify users or correlate credential
presentations. Issuers, users, and verifiers may be
arbitrarily malicious, and may collude to try to undermine the privacy of honest
users, to cause honest verifiers to accept invalid presentations, or to forge
credentials.  We consider static corruption throughout: the adversary fixes the
set of corrupt participants before any protocol interaction takes place.

An aggregate anonymous credential scheme $\AAC$ consists of the following
algorithms.

\begin{description}

  \item[$\crs \sample \AACgen(\secparam, n)$ :] The setup algorithm, on input
security parameter $\secparam$ and number of issuers $n$, generates the system
public parameters $\crs$, which include the public keys of the registration
authority and of the $n$ issuers.

  \item[$(\sk,\regcred) \sample \AACregister(\crs,\rask)$ :] The registration
algorithm, run between a user and the registration authority, on input public
parameters $\crs$ and the registration authority's secret key $\rask$, outputs
user secret key $\sk$ and registration credential $\regcred$.

  \item[$\cred \sample \AACissue(\crs,\isk{i},\sk,\regcred,\atv{i})$ :] The
issuance algorithm, run between a user and issuer $i$, on input public
parameters $\crs$, issuer secret key $\isk{i}$, user secret key $\sk$,
registration credential $\regcred$, and value $\atv{i}$ of attribute
$\att{i}$, outputs attribute credential $\cred$.

  \item[$\aggcred \leftarrow \AACaggregate(\crs,(\cred_i)_{i \in \ISSUERS})$ :]
The aggregation algorithm, on input public parameters $\crs$ and credentials
$(\cred_i)_{i \in \ISSUERS}$ for a set of issuers $\ISSUERS \subseteq [1,n]$,
outputs aggregate credential $\aggcred$.

  \item[$\algoproof \sample \AACpresent(\crs,\sk,\regcred,\aggcred,\msg)$ :]
The presentation algorithm, on input public parameters $\crs$, user secret key
$\sk$, registration credential $\regcred$, aggregate credential $\aggcred$, and
message $\msg$, chosen by the verifier and bound to the presentation, outputs
presentation $\algoproof$.

  \item[$\{0,1\} \leftarrow \AACverify(\crs,\{\atv{i}\}_{i \in
\ISSUERS},\msg,\algoproof)$ :] The verification algorithm, on input public
parameters $\crs$, disclosed attribute values $\{\atv{i}\}_{i \in \ISSUERS}$
for a set of issuers $\ISSUERS \subseteq [1,n]$, message $\msg$, and
presentation proof $\algoproof$, returns 0 or 1.

\end{description}

\subsection{Tag-based aggregate signature schemes}

We define a tag-based aggregate signature (ATS) scheme as one that allows
signatures computed over the same tag to be aggregated into a compact signature.
A tag is a string that is signed together with each message, and only signatures
under the same tag can be aggregated. With $n$ signers and a set $\ISSUERS
\subseteq [1,n]$, a tag-based aggregate signature scheme is defined by
the following algorithms:

\begin{description}

  \item[$(\sk, \pk) \sample \ATSgen(\secparam)$ :] The generation algorithm, on
input security parameter $\secparam$, outputs key pair $(\sk,\pk)$.

  \item[$\signature \sample \ATSsign(\sk, \Tag, \msg)$ :] The signing algorithm,
on input secret key $\sk$, tag $\Tag \in \{0,1\}^*$, and message $\msg \in
\{0,1\}^*$, outputs signature $\signature$.

  \item[$ \signature' \sample \ATSrandomise(\Tag, \signature)$ :] The
randomisation algorithm, on input tag $\Tag$ and signature $\signature$, outputs
randomised signature $\signature'$.

  \item[$ \{0,1\} \leftarrow \ATSverify(\pk, \Tag, \msg, \signature)$ :] The
verification algorithm, on input public key $\pk$, tag $\Tag$, message $\msg$,
and signature $\signature$, returns 0 or 1.

  \item[$\signature \leftarrow \ATSaggregate((\signature_i)_{i \in \ISSUERS})$
:] The aggregation algorithm, for set $\ISSUERS \subseteq [1, n]$, on input
signatures $(\signature_i)_{i \in \ISSUERS}$, all under the same tag $\Tag$,
outputs aggregate signature $\signature$.

  \item[$\{0,1\} \leftarrow \ATSverifyagg((\pk_i)_{i\in \ISSUERS}, \Tag,
(\msg_i)_{i\in \ISSUERS}, \signature)$ :] The aggregate verification algorithm,
for set $\ISSUERS \subseteq [1, n]$, on input public keys $(\pk_i)_{i\in
\ISSUERS}$, tag $\Tag$, messages $(\msg_i)_{i\in \ISSUERS}$, and aggregate
signature $\signature$, returns 0 or 1.

\end{description}

\paragraph{Correctness.} An $\ATS$ scheme satisfies correctness if the
following holds for all security parameters $\secparam$, all $n \geq 1$, all
tags $\Tag \in \{0,1\}^*$, and all messages $\msg_i \in \{0,1\}^*$. First,
the scheme satisfies the usual correctness definition for digital signature
schemes, so for every key pair $(\sk,\pk)$ generated by $\ATSgen(\secparam)$ and
every tag and message $\Tag, \msg$, the following holds:
%
\begin{equation*}
%
  \ATSverify(\pk, \Tag, \msg, \ATSsign(\sk, \Tag, \msg)) = 1.
%
\end{equation*}
%

We also have correctness of aggregated signatures.  For set $\ISSUERS
\subseteq [1, n]$, for $i \in \ISSUERS$, an aggregate signature is generated as
follows:
%
\begin{gather*}
%
  (\sk_i, \pk_i) \sample \ATSgen(\secparam), \quad \signature_i \sample
\ATSsign(\sk_i, \Tag, \msg_i), \\ \signature \gets
\ATSaggregate((\signature_i)_{i \in \ISSUERS}).
    %
\end{gather*}
%
Then aggregate correctness holds if we have the following:
  %
\begin{gather*}
%
  \text{If } \ \forall i \in \ISSUERS \ : \ \ATSverify(\pk_i, \Tag, \msg_i,
\signature_i) = 1 \\ \text{then } \ \ATSverifyagg((\pk_i)_{i \in \ISSUERS},
\Tag, (\msg_i)_{i \in \ISSUERS}, \signature) = 1.
    %
\end{gather*}
%
Randomisation also preserves validity, meaning that the following holds:
  %
\begin{gather*}
%
  \text{If } \ \ATSverify(\pk, \Tag, \msg, \signature) = 1, \\ \text{then }
\forall \ \signature' \ : \  \signature' \sample \ATSrandomise(\Tag,
\signature), \\ \ATSverify(\pk, \Tag, \msg, \signature') = 1.
    %
\end{gather*}
%

\paragraph{Existential unforgeability under tag-based aggregation (EUF-TBA).}
The advantage of a PPT
adversary $\adv$ with access to corruption oracle $\corrupto$ and signing
oracle $\signo$ is defined as follows:
%
\begin{equation*}
%
  \Adv_{\ATS, \adv}^{\mathsf{EUF\text{-}TBA}}(\secparam) = \Pr[\adv \text{
wins}],
%
\end{equation*}
%
where the experiment and its oracles $\corrupto$ and $\signo$ are given in
\Cref{exp:ats-euf}. An $\ATS$
scheme is existentially unforgeable under tag-based aggregation if this
advantage is negligible for every PPT adversary $\adv$.

\begin{experimentbox}{Tag-based aggregate signature schemes: existential unforgeability under tag-based aggregation}{ats-euf}

  \begin{enumerate}

    \item The challenger generates $(\sk_i, \pk_i) \sample \ATSgen(\secparam)$ for
each $i \in [1,n]$, and sends $(\pk_1, \ldots, \pk_n)$ to $\adv$.

    \item The adversary is given oracle access to $\corrupto$, which on input $i
\in [1,n]$ adds $\pk_i$ to $\Cor$ and returns $\sk_i$.

    \item The adversary is given oracle access to $\signo$, which on input $(i,
\Tag, \msg)$ with $i \in [1,n]$ adds $(\pk_i, \Tag, \msg)$ to $\Query$ and
returns $\signature \sample \ATSsign(\sk_i, \Tag, \msg)$.

    \item The adversary outputs $(\ISSUERS \subseteq [1, n], (\pk'_i)_{i \in
\ISSUERS}, \Tag, (\msg_i)_{i \in \ISSUERS}, \signature)$, where $\pk'_i =
\pk_i$ for every $i \in \ISSUERS$ with $\pk_i \notin \Cor$, and otherwise the
adversary also outputs $\sk'_i$ such that $(\sk'_i, \pk'_i)$ is a key pair of
$\ATSgen$.

    \item Challenger computes $b \gets \ATSverifyagg\bigl((\pk'_i)_{i \in
\ISSUERS}, \Tag, (\msg_i)_{i \in \ISSUERS}, \signature\bigr)$.

  \end{enumerate}

  $\adv$ wins the experiment if $b = 1$ and $\exists \ i \in \ISSUERS : \pk_i
\notin \Cor \wedge (\pk_i, \Tag, \msg_i) \notin \Query$.

\end{experimentbox}

Corrupted keys are certified, so the adversary outputs the secret key of every
corrupted key it uses.

\subsection{The ideal functionality for aggregate anonymous credentials}

We define security for the aggregate anonymous credentials protocol through the
ideal functionality $\idfu{\AAC}$, in the sense of \Cref{sec:prelims-security}.
The formal description is given in~\Cref{fig:agg-creds-faac}. We summarise each
interface below.

The functionality keeps three tables. The first, $\creds[\user]$, lists for
each issuer $i$ the set of values of $\att{i}$ issued to user $\user$. This set
is $\varnothing$ for an issuer that has issued nothing to $\user$, and the whole
entry is $\bot$ while $\user$ is unregistered. The second,
$\aggcreds[\signature]$, maps an aggregate credential $\signature$ to its user
and values. The third, $\proven[\vec{\atv{}}, \msg, \algoproof]$, records what
is known about a presentation. It is $1$ once the presentation is known to be
valid, $0$ once it is known to be invalid, and $\bot$ while neither is known. We
write $T[x] \rightarrow y$ when entry $x$ of table $T$ is $y$, and $T[x]
\not\rightarrow \bot$ when it is set. Issuer $i$ is $\issuer{i}$, and a verifier
is $\algo{V}$.

\paragraph{Register.} Registration lets a user join the system. If the
registration authority accepts and the user was never registered before, it
succeeds.

\paragraph{Issue.} A user obtains a credential certifying value $\atv{}$ of
attribute $\att{}$. If the authorised issuer accepts, the user
receives a credential usable in later presentations. Honest issuers only
issue to registered users. An issuer may certify several values of its attribute
for the same user, and a presentation discloses at most one of them. The
functionality records every value issued, and security does not depend on how
many credentials a user holds. $\simu$ learns the user's identity on
every call, and also the attribute value on issue calls.

\paragraph{Aggregate.} The user's credentials for attribute values
$\vec{\atv{}}$ are combined into a form usable for later presentation.
Aggregation is local, so $\simu$ learns nothing. Unregistered users cannot
aggregate, and registered users aggregate exactly the values they hold. Success
results in a random string $\signature$ representing the aggregate credential.

\paragraph{Present.} To produce a presentation of $\vec{\atv{}}$, the user
submits $(\vec{\atv{}}, \signature, \msg)$. Here
$\signature$ aggregates the credentials for $\vec{\atv{}}$, and $\msg$ is a fresh
message from the verifier to which the presentation is bound.
$\idfu{\AAC}$ accepts if the user previously called aggregate on $\vec{\atv{}}$,
or otherwise if the user actually holds $\vec{\atv{}}$. This means a corrupt
user who skips aggregate can still succeed if they hold the credentials.
Unregistered users cannot present. On
acceptance, $\idfu{\AAC}$ generates a random string $\algoproof$, so the
presentation reveals nothing about the user beyond the disclosed values.

\paragraph{Verify.} A verifier checks a presentation
$(\vec{\atv{}}, \msg, \algoproof)$. If it was already produced by a valid
present call, $\idfu{\AAC}$ returns true. If already flagged invalid, it returns
false.  Otherwise, $\idfu{\AAC}$ checks whether the presentation could have been
forged by corrupt parties: whether corrupt users hold all disclosed values, or
failing that, whether corrupt issuers could produce valid credentials for the
remaining values. If either check succeeds, $\idfu{\AAC}$ defers to $\simu$ and
returns its answer. Otherwise, it records the presentation as invalid and
returns false. The first case gives completeness, since every honest
presentation verifies. The last gives unforgeability, since a presentation that
corrupt parties could not have produced is rejected.

    \begin{framed}

      \setlist[enumerate]{ topsep=2pt, partopsep=0pt, parsep=0pt, itemsep=1pt }

	\fuinterface{\AACregister} On input register from the user, send $(register,
user)$ to the registration authority and wait for its confirmation.  On
receiving the confirmation, do the following:

	\begin{enumerate}

	  \item If $\creds[\user] \not\rightarrow \bot$, then stop.

	  \item Send $(\AACregister, \user)$ to $\simu$ and wait for the
go-ahead. On receiving the go-ahead, continue.

	  \item Set $\creds[\user] \leftarrow [\varnothing, \ldots,
\varnothing]$ for length $n$.

	  \item Return $\AACregisterend$ to the user.

	\end{enumerate}

	\fuinterface{\AACissue} On input $(\AACissue, \att{i},\atv{i})$ from the user,
with $\atv{i}$ being value of attribute $\att{i}$, send $(\AACissue, \user,
\atv{i})$ to issuer $\issuer{i}$ associated with $\att{i}$ and wait for its
confirmation.  On confirmation, do:

	\begin{enumerate}

	  \item If issuer $\issuer{i}$ is honest, then:

	    \begin{enumerate}

	      \item If $\creds[\user] \rightarrow \bot$, then stop.

	      \item Send $(\AACissue, \user, \att{i}, \atv{i})$ to $\simu$ and
wait for the go-ahead. On receiving it, continue.

	    \end{enumerate}

	  \item Otherwise:

	    \begin{enumerate}

	      \item Send $(\AACissue, \user, \att{i}, \atv{i})$ to $\simu$, and
wait for the go-ahead. On receiving it, continue.

	    \end{enumerate}

	  \item Set $\creds[\user][i] \leftarrow \creds[\user][i] \cup
\{\atv{i}\}$.

	  \item Return $(\AACissueend, \att{i}, \atv{i})$ to the user and
$(\AACissueend, \user, \atv{i})$ to $\issuer{i}$.

	\end{enumerate}

	\fuinterface{\AACaggregate} On input $(\AACaggregate, \vec{\atv{}})$ from the
user, where $\vec{\atv{}} = (\atv{1}, \ldots, \atv{n})$ and $\atv{i} \in
\{\varnothing\} \cup \{0, 1\}^*$, do:

	\begin{enumerate}

	  \item If $\creds[\user] = \bot$, then stop.

	  \item If $\exists i \in [1, n]: \atv{i} \neq \varnothing \land \atv{i}
\notin \creds[\user][i]$, then stop. Otherwise:

	    \begin{enumerate}

	      \item Randomly generate $\signature$.

	      \item Set $\aggcreds[\signature] \leftarrow \langle \user,
\vec{\atv{}}\rangle$.

	      \item Return $(\AACaggregateend, \signature)$ to the user.

	    \end{enumerate}

	\end{enumerate}

	\fuinterface{\AACpresent} On input $(\AACpresent, \vec{\atv{}}, \signature,
\msg)$ from the user, where $\vec{\atv{}} = (\atv{1}, \ldots, \atv{n})$, with
$\atv{i} = \varnothing$ if $\atv{i}$ is not disclosed, do:

	\begin{enumerate}

	  \item If $\aggcreds[\signature] \rightarrow \langle \user,
\vec{\atv{}}\rangle$, then:

	    \begin{enumerate}

	      \item Randomly generate $\algoproof$.  \item Set
$\proven[\vec{\atv{}}, \msg, \algoproof] \leftarrow 1$.

	      \item Return $(\AACpresentend, \algoproof)$ to the user and
$(\AACpresentend, \vec{\atv{}}, \msg, \allowbreak \algoproof)$ to $\simu$.

	    \end{enumerate}

	  \item Otherwise, if $\creds[\user] = \bot$, then stop.

	  \item If $\exists i \in [1, n]: \atv{i} \neq \varnothing \land \atv{i}
\notin \creds[\user][i]$, then stop. Otherwise:

	    \begin{enumerate}

	      \item Randomly generate $\algoproof$.

	      \item Set $\proven[\vec{\atv{}}, \msg, \algoproof] \leftarrow 1$.

	      \item Return $(\AACpresentend,  \algoproof)$ to the user and
$(\AACpresentend, \vec{\atv{}}, \msg, \allowbreak \algoproof)$ to $\simu$.

	    \end{enumerate}

	\end{enumerate}

	\fuinterface{\AACverify} On receiving message $(\AACverify, \vec{\atv{}},
\allowbreak \msg, \algoproof)$ from a verifier, do:

	\begin{enumerate}

	  \item If $\proven[\vec{\atv{}}, \msg, \algoproof] \rightarrow 0$, then
return $(\AACverifyend, 0)$ to $\algo{V}$.

	  \item Otherwise, if $\proven[\vec{\atv{}}, \msg, \algoproof]
\rightarrow 1$, then:

	    \begin{enumerate}

	      \item Return $(\AACverifyend, 1)$ to the verifier.

	      \item Send $(\AACverifyend, \algo{V}, \vec{\atv{}}, \msg,
\algoproof, \allowbreak 1)$ to $\simu$.

	    \end{enumerate}

	  \item Otherwise, if $\proven[\vec{\atv{}}, \msg, \algoproof]
\rightarrow \bot$, then:

	    \begin{enumerate}

	      \item If there exists a dishonest user $\user'$ such that for every $i$ with $\atv{i} \neq \varnothing$, either $\atv{i} \in
\creds[\user'][i]$ or issuer $\issuer{i}$ is dishonest, then query $\simu$ with
$(\AACverify, \vec{\atv{}}, \msg,\algoproof)$. On receiving the response
$(\AACverifyend, b)$, do:

		\begin{enumerate}

		  \item Set $\proven[\vec{\atv{}}, \msg, \algoproof] \leftarrow
b$.

		  \item Return $(\AACverifyend, b)$ to $\algo{V}$.

		\end{enumerate}

	      \item Otherwise:

		\begin{enumerate}

		  \item Set $\proven[\vec{\atv{}}, \msg, \algoproof] \leftarrow
0$.

		  \item Return $(\AACverifyend, 0)$ to $\algo{V}$.

		  \item Send $(\AACverifyend, \algo{V}, \vec{\atv{}}, \msg,
\allowbreak \algoproof, \allowbreak 0)$ to $\simu$.

		\end{enumerate}

	    \end{enumerate}

	\end{enumerate}

    \end{framed}

\captionof{figure}[Ideal functionality for aggregate anonymous credentials]{Ideal functionality for aggregate anonymous
credentials.}\label{fig:agg-creds-faac}
